% from ega4-9-fr.tex
\oldpage[IV-3]{80}

irréductibles (resp. géométriquement connexes) et que $(X_\eta)_{\mathrm{red}}$
est géométriquement réduit.

Supposons d'abord que $\eta\in S-E$ pour une des propriétés (i) à (iv). Si
$X_\eta$ n'est pas géométriquement irréductible (resp. géométriquement connexe),
il n'est pas irréductible (resp. connexe) en vertu de l'hypothèse précédente,
donc il en est de même de $X_s$ pour $s$ dans un voisinage de $\eta$ (9.7.1), et
\emph{a fortiori} dans ce voisinage $X_s$ n'est pas géométriquement irréductible
(resp. géométriquement connexe). D'autre part, si $X_\eta$ n'est pas
géométriquement réduit, il n'est pas réduit (sans quoi il serait égal à
$(X_\eta)_{\mathrm{red}}$ qui est géométriquement réduit par hypothèse)~; donc,
dans un voisinage de $\eta$, $X_s$ n'est pas réduit (9.7.2) et \emph{a fortiori}
n'est pas géométriquement réduit. Enfin, si $X_\eta$ n'est pas géométriquement
intègre, ou bien il n'est pas réduit, et alors on vient de voir que $X_s$ n'est
pas réduit (ni \emph{a fortiori} intègre) dans un voisinage de $\eta$~; ou bien
$X_\eta$ est réduit (donc géométriquement réduit par hypothèse), et alors il
n'est pas géométriquement irréductible, et on a vu ci-dessus qu'il en est alors
de même de $X_s$ pour $s$ voisin de $\eta$~; \emph{a fortiori} $X_s$ n'est pas
géométriquement intègre pour ces valeurs de $s$.

Nous allons donc désormais supposer que $\eta\in E$ et examiner séparément
chacune des propriétés considérées.

1$^\circ$ Supposons que $X_\eta$ soit géométriquement intègre. Soit $L$ le corps
des fonctions rationnelles sur $X_\eta$~; l'hypothèse sur $X_\eta$ implique que
$L$ est extension séparable de $K$ (4.6.3), donc extension séparable finie d'une
extension pure $K(\mathrm T_1,\ldots,\mathrm T_n)$ ($\mathrm T_i$
indéterminées)~; il y a par suite un élément $x\in L$, entier sur l'anneau
$K[\mathrm T_1,\ldots,\mathrm T_n]$, et tel que
$L=K(\mathrm T_1,\ldots,\mathrm T_n)(x)$~; soit
$G\in K[\mathrm T_1,\ldots,\mathrm T_n,\mathrm T_{n+1}]$ son polynôme minimal.
Il existe un élément $g\neq0$ de $A$ tel que tous les coefficients $\neq0$ de
$G$ (qui appartiennent à $K$) soient dans l'anneau $A_g$~; remplaçant $A$ par
$A_g$ (ce qui revient à remplacer $S$ par un voisinage de $\eta$), on peut donc
supposer que $G$ a ses coefficients dans $A$~; désignant par $F$ le polynôme
$G$ considéré comme élément de
$A[\mathrm T_1,\ldots,\mathrm T_n,\mathrm T_{n+1}]$, on a donc
$G=F_{(\boldsymbol{k}(\eta))}$. Posons
$Y=\operatorname{Spec}(A[\mathrm T_1,\ldots,\mathrm T_{n+1}]/(F))$, de sorte
que
$Y_\eta=\operatorname{Spec}(K[\mathrm T_1,\ldots,\mathrm T_{n+1}]/(G))$~;
$Y_\eta$ est un schéma intègre ayant $L$ pour corps des fonctions rationnelles.
Comme $X_\eta$ et $Y_\eta$ sont noethériens, il existe un ouvert non vide
$V\subset Y_\eta$ et une immersion ouverte $v:V\to X_\eta$ (nécessairement
dominante) (\textbf{I}, 6.5.1, (ii) et 6.5.4, (ii)). Soit $W$ un ouvert de $Y$
tel que $W\cap Y_\eta=V$~; appliquant (8.8.2, (i)) et (8.10.5, (iii)) suivant
la méthode de (8.1.2, \emph{a)}), on voit qu'en remplaçant au besoin $S$ par un
voisinage de $\eta$, on peut supposer que $v=w_\eta$, où $w:W\to X$ est une
immersion ouverte.

Cela étant, on a vu (4.6.3) que le critère pour qu'un préschéma algébrique
intègre soit géométriquement intègre ne dépend que de son corps de fonctions
rationnelles~; comme $X_\eta$ est géométriquement intègre par hypothèse, il en
est de même de $Y_\eta$, et la définition de $G$ entraîne donc que ce polynôme
est géométriquement irréductible (9.7.4). Appliquant (9.7.5), on voit qu'il y a
un voisinage $U$ de $\eta$ dans $S$ tel que $F_{(\boldsymbol{k}(s))}$ soit
géométriquement irréductible pour tout $s\in U$, donc $Y_s$ est géométriquement
intègre pour $s\in U$ (9.7.4)~; en outre on peut supposer que pour $s\in U$,
$W_s$ n'est pas vide (9.5.1), et par suite est géométriquement intègre (4.6.3)~;
enfin, on peut supposer aussi que

